\section{Introducción}
\label{sec:intro}

\subsection{Problema y contexto}

La palma aceitera es un cultivo perenne cuya rentabilidad depende de decisiones que se toman en
ventanas de tiempo estrechas y sobre información que hoy no se registra. El proyecto se desarrolla
con una organización productora del cantón El Empalme, provincia del Guayas, que explota
aproximadamente cien hectáreas distribuidas en seis lotes y emplea a sesenta y dos personas.

El trabajo de elicitación de las unidades anteriores documentó tres problemas que el sistema debe
atender y que conviene enunciar antes que cualquier consideración técnica, porque son los que
justifican el proyecto:

\begin{enumerate}[leftmargin=*, nosep]
  \item \textbf{La detección fitosanitaria depende de una visita quincenal.} Entre visitas de la
        asesoría técnica, una afectación avanza sin que nadie la registre. El personal de campo la
        observa pero no dispone de un medio para clasificarla ni para dejar constancia.
  \item \textbf{La penalización por fruta verde se decide sobre un juicio visual.} La planta
        extractora aplica un descuento cuando el porcentaje de fruta verde supera un umbral, y la
        decisión de cortar se toma con una estimación hecha a ojo en el momento del corte.
  \item \textbf{La liquidación semanal se calcula a mano.} La evidencia \id{EV-13} documenta que el
        cálculo de la remuneración del personal consume una jornada administrativa completa cada
        semana y se realiza sobre anotaciones en papel.
\end{enumerate}

Ninguno de los tres es un problema de software. Son problemas de operación agrícola cuya solución
pasa por registrar información que hoy se pierde, y esa distinción es la que gobierna todo el
documento: el SIMPA no automatiza porque se pueda, sino donde la falta de registro tiene una
consecuencia económica medible sobre la organización o sobre las personas que trabajan en ella.

\subsection{El sistema}

El SIMPA ---Sistema Inteligente de Mantenimiento de Palma Africana--- es una aplicación móvil con
panel web de administración que cubre siete bloques funcionales: gestión de la estructura productiva,
planificación y seguimiento de labores, registro operativo de campo, polinización asistida,
diagnóstico asistido por inteligencia artificial, calidad de la fruta y relación con la extractora, y
reportes, estimación e histórico.

La versión final del ERS especifica 42 requisitos funcionales, 19 no funcionales, 10 restricciones de
diseño y 8 requisitos legales derivados de la Ley Orgánica de Protección de Datos Personales, más 18
requisitos específicos de los dos componentes de inteligencia artificial que esta unidad incorpora.

\subsection{Objetivos de esta práctica}

\noindent\textbf{Objetivo general.} Integrar, auditar y defender el proyecto: aplicar métricas de
calidad al ERS final, cerrar la trazabilidad extremo a extremo, especificar los requisitos de los
componentes de inteligencia artificial y sostener las decisiones de ingeniería ante el tribunal.

\noindent\textbf{Objetivos específicos.}

\begin{enumerate}[leftmargin=*, nosep]
  \item Construir un instrumento de auditoría con seis métricas calculadas sobre conteos verificables
        de las fuentes del propio documento.
  \item Corregir toda métrica por debajo de su valor de referencia y volver a medir, presentando el
        par antes/después.
  \item Cerrar la trazabilidad extremo a extremo, incorporando los eslabones de proceso, estado y
        criterio de aceptación que la cadena anterior no tenía.
  \item Especificar dos componentes de inteligencia artificial con sus requisitos de rendimiento,
        equidad y explicabilidad, su clasificación de riesgo y su plan de monitoreo.
  \item Consolidar la versión final del ERS y documentar el aporte individual verificable.
\end{enumerate}

\subsection{Estructura del documento}

La \S\ref{sec:metodologia} describe el proceso de ingeniería de requisitos seguido a lo largo de las
cinco unidades y justifica las decisiones metodológicas. La \S\ref{sec:ers} presenta la estructura
del ERS final y sus requisitos clave. La \S\ref{sec:uml} recorre los modelos con lectura
interpretativa. La \S\ref{sec:validacion} expone la inspección formal, sus defectos y la
re-inspección de esta unidad. La \S\ref{sec:gestion} trata la línea base, el control de cambios y la
matriz de trazabilidad final. La \S\ref{sec:ia} especifica los componentes de inteligencia
artificial. La \S\ref{sec:metricas} presenta la auditoría de calidad con la aritmética de cada
métrica. La \S\ref{sec:retro} recoge la retrospectiva y la \S\ref{sec:conclusiones} cierra con cinco
conclusiones, una por unidad.

